\documentclass[11pt]{article}
\usepackage[a4paper,margin=1in]{geometry}
\usepackage{amsmath,amssymb,amsthm}
\usepackage{booktabs}
\usepackage{enumitem}
\usepackage{hyperref}
\usepackage{mathtools}
\usepackage{bm}

\setlist{nosep,leftmargin=*}
\setlength{\parskip}{3pt}
\setlength{\parindent}{0pt}

\title{Top Three Algorithmic Improvement Proposals for Vamana ANN\\
\large Extracted from Sections 2, 3, and 4}
\author{Prepared for the GraphANN repository}
\date{\today}

\begin{document}
\maketitle

\begin{abstract}
This note summarizes three high-impact improvements for a baseline Vamana ANN index, with concise implementation guidance, expected impact, and mathematical rationale.
\end{abstract}

\section{Improvement 1: Two-Pass Vamana with Medoid Start and $\alpha$ Scheduling}

\textbf{Implementation.}
\begin{enumerate}
    \item Replace random entry node with a medoid-like center:
    \begin{equation}
    \mu = \frac{1}{n}\sum_{i=1}^n x_i,\qquad s = \arg\min_i \lVert x_i - \mu \rVert_2^2.
    \end{equation}
    \item Start from a lightly random-initialized graph, then run two insertion passes: first with $\alpha_1=1$, second with $\alpha_2>1$ (e.g., $1.1$--$1.4$).
    \item Keep bounded degree via target $R$ and overflow threshold $\gamma R$.
\end{enumerate}

\textbf{Why better.} Medoid entry reduces average routing distance; two-pass construction reduces insertion-order bias; $\alpha_1=1$ preserves local structure while $\alpha_2>1$ restores long-range edges, typically lowering hop count.

Define expected hop count to an acceptable neighborhood of true nearest neighbors as $H(q;G,s)$. Under fixed degree budget, we seek:
\begin{equation}
\min_G\; \mathbb{E}_q[H(q;G,s)]\quad\text{s.t.}\quad \deg^+(v)\le R.
\end{equation}

Two-pass Vamana is a continuation in pruning aggressiveness:
\begin{equation}
G^{(1)} = \mathcal{P}_{\alpha=1}(G_0),\qquad G^{(2)} = \mathcal{P}_{\alpha=\alpha_2}(G^{(1)}),\quad \alpha_2>1.
\end{equation}

\textbf{Cost.} Build is usually slower ($\approx1.7\times$--$2.0\times$), but target recall can often be achieved with smaller query-time $L$.

\section{Improvement 2: Candidate Set for Prune from Visited Set (Not Only Final Frontier)}

\textbf{Implementation.}
During insertion of point $p$, greedy search should expose two sets:
\begin{itemize}
    \item Final bounded frontier $S_L(p)$ (current behavior).
    \item Full visited set $V_p$.
\end{itemize}

Use prune candidates
\begin{equation}
C_p = V_p \cup N_{\text{out}}(p)
\end{equation}
instead of only frontier nodes, then robust-prune $C_p$ to at most $R$ edges.

\textbf{Why better.} The final frontier is lossy; visited nodes often include bridge vertices removed by truncation. A richer candidate pool increases diversity and improves navigability.

Let $\mathcal{F}(A)$ be prune utility under proximity/diversity constraints. Current selection solves
\begin{equation}
\max_{A\subseteq S_L(p),\,|A|\le R} \mathcal{F}(A).
\end{equation}

Improved selection solves
\begin{equation}
\max_{A\subseteq (V_p\cup N_{\text{out}}(p)),\,|A|\le R} \mathcal{F}(A).
\end{equation}

Since $S_L(p)\subseteq V_p\cup N_{\text{out}}(p)$, the best achievable objective cannot decrease. \textbf{Cost:} more pruning work, often offset by lower query-time beam for equal recall.

\section{Improvement 3: Locally Adaptive $\alpha_i$ and Degree Budget $R_i$}

\textbf{Implementation.}
\begin{enumerate}
    \item Estimate local intrinsic dimensionality (LID) per point $i$ from local distance statistics:
    \begin{equation}
    \widehat{m}_i = \left(-\frac{1}{k-1}\sum_{j=1}^{k-1}\log\frac{r_j(i)}{r_k(i)}\right)^{-1},
    \end{equation}
    where $r_j(i)$ is the $j$-th nearest local sample distance.
    \item Set node-specific pruning aggressiveness and degree:
    \begin{equation}
    \alpha_i = \mathrm{clip}\left(1+\frac{c_\alpha}{\sqrt{\widehat{m}_i}},\;\alpha_{\min},\alpha_{\max}\right),
    \end{equation}
    \begin{equation}
    R_i = \mathrm{clip}\left(R_0 + c_R\widehat{m}_i,\;R_{\min},R_{\max}\right).
    \end{equation}
    \item Use per-node values in robust prune and overflow checks (threshold $\gamma R_i$).
    \item Enforce global memory budget:
    \begin{equation}
    \sum_{i=1}^n R_i \le nR_{\text{budget}}.
    \end{equation}
\end{enumerate}

\textbf{Why better.} Real data is locally heterogeneous; fixed $\alpha,R$ misallocate edges. Adaptive budgets allocate more connectivity to hard regions and reduce waste in easy regions.

Global policy:
\begin{equation}
\min_G \sum_i \mathcal{L}_i(G;\alpha,R).
\end{equation}
Adaptive policy:
\begin{equation}
\min_G \sum_i \mathcal{L}_i(G;\alpha_i,R_i)
\quad\text{s.t.}\quad
\sum_i R_i \le nR_{\text{budget}}.
\end{equation}

This is constrained resource allocation; LID acts as a local complexity proxy. \textbf{Cost:} light metadata/estimation overhead, with potential recall-latency gains at fixed average memory.

\end{document}
